\documentclass[11pt,a4paper]{article}
\usepackage[T1]{fontenc}
\usepackage[utf8]{inputenc}
\usepackage[ngerman]{babel}
\usepackage[margin=2.5cm]{geometry}
\usepackage[hidelinks]{hyperref}
\setlength{\parindent}{0pt}
\setlength{\emergencystretch}{2em}
\begin{document}
\begin{center}
  {\Large\textbf{Webtechnologien -- Übungsblatt 1}}
\end{center}
\section*{Aufgabe 2: Apache}
\textbf{Aufgabe:} Apache einrichten, Konfiguration und localhost
erklären, die Startseite ändern und den Ordner \texttt{mytest} testen.

\medskip
\textbf{Ergebnis:} Apache 2.4.58 ist über XAMPP installiert und läuft.
Die Startseite unter \url{http://localhost/} zeigte zuerst
\glqq It works\grqq. Nach Änderung der \texttt{index.html}
steht dort \glqq Apache läuft!\grqq.
Bei XAMPP wurde \texttt{index.html} über \texttt{DirectoryIndex}
vor \texttt{index.php} gesetzt.

\begin{description}
  \item[httpd.conf:] Hauptkonfiguration für Ports, Verzeichnisse, Module
    und Zugriffsregeln; hier unter \path{C:/xampp/apache/conf/httpd.conf}.
  \item[ServerRoot:] Basisordner der Apache-Installation:
    \path{C:/xampp/apache}.
  \item[DocumentRoot:] Ordner der ausgelieferten Webseiten:
    \path{C:/xampp/htdocs}.
  \item[Virtual Host:] Ermöglicht mehrere Webseiten mit eigener
    Konfiguration auf einem Apache-Server.
  \item[127.0.0.1 und localhost:] Loopback-Adresse und lokaler Hostname.
    Damit wird der eigene Rechner angesprochen.
  \item[index.html:] Standard-Startseite eines Verzeichnisses,
    sofern als \texttt{DirectoryIndex} konfiguriert.
\end{description}

\textbf{mytest:} Der Ordner \path{C:/xampp/htdocs/mytest} enthält
zwei Textdateien und einen Unterordner mit einer weiteren Datei.
Beim Aufruf von
\url{http://localhost/mytest/} erscheint \glqq Index of /mytest\grqq{}
mit den Dateien und dem Unterordner. Es gibt keine Indexdatei;
\texttt{Options Indexes} erlaubt die Verzeichnisliste.

\medskip
\textbf{Test:} Startseite und Verzeichnisliste liefern \texttt{200 OK}.

\medskip
\textbf{Quelle:} \href{https://httpd.apache.org/docs/2.4/mod/core.html}{Apache-Dokumentation}.
\medskip
{\small \textbf{Hilfsmittel:} Codex für Erstellung, Prüfung und Dokumentation.}
\end{document}
